% Section 2.7, from lectures/L02/appendix/a_larger_networks.md (A.7) and
% lectures/L02/appendix/c_testbenches.md (C.1-C.6).

\section{Verifying your VHDL with a testbench}
\label{c2:sec:testbenches}

Every VHDL module you write for an exercise comes with a \textbf{self-checking testbench}: a small
VHDL file that drives your inputs, checks your outputs and says which case failed.
\Chapref{c1:ch}'s exercises already handed such things out, and \secref{c1:sec:module} gave you the
three commands that run one.

You are not asked to write the testbench that verifies a module you build; every single one of them
is handed out. Running them is enough to verify every exercise on your own laptop, with no FPGA
board. (A couple of project exercises from \chapref{c13:ch} onwards invite a short test rig of your
own, but nothing you are assessed on rests on a testbench you wrote yourself.)

This section is the book's introduction, and every later chapter's exercises refer back here.
Appendix~\ref{app:sim} is the fuller version, written for the group project; come back here for the
introduction and go there for the details.

\subsection{What you need}
\label{c2:sub:need}
\textbf{GHDL}, the open-source VHDL analyser and simulator. Check with \code{ghdl --version};
install on WSL or Ubuntu with \code{sudo apt -y install ghdl}.

Every command below uses \code{--std=93}, the VHDL standard the book targets.

\subsection{What a testbench is}
\label{c2:sub:what}
A testbench is itself a VHDL module, but a special one:
\begin{itemize}
  \item Its entity has \textbf{no ports}: nothing connects to the outside world, since it is the
    top of the simulation.
  \item Its architecture \textbf{instantiates your design}, the ``device under test'', labelled
    \code{dut}, and drives its inputs.
  \item It \textbf{checks} every output with an \code{assert}. If an output is wrong the assertion
    fires with a message and the simulation stops.
\end{itemize}

So instead of you reading a waveform diagram to decide whether the output is right, the testbench
decides it, and speaks up only when something is wrong.

\subsection{Running a testbench: analyse, elaborate, run}
\label{c2:sub:run}
Three steps. Pass \code{--std=93} to all three, and give the \textbf{testbench's entity name}, not
a file name, to the last two:

\begin{shell}
ghdl -a --std=93 <module>.vhd <module>_tb.vhd        # analyze both files (module first)
ghdl -e --std=93 <module>_tb                         # elaborate the testbench
ghdl -r --std=93 <module>_tb --assert-level=error --stop-time=10ms   # run it
\end{shell}

\begin{itemize}
  \item \code{--assert-level=error} makes a failed check stop the run and return a
    \textbf{non-zero exit code}, a real pass/fail that you can check with \code{echo $?}. The
    testbenches in this book's first part raise failures with \code{severity failure}, which stops
    the run by itself, so you get a real pass/fail even without the flag. Keep it on anyway: it is
    the habit that also works with testbenches reporting at a lower severity level, and the group
    project's testbenches are exactly those (appendix~\ref{app:sim}).
  \item \code{--stop-time=10ms} caps how long the simulation may run. A testbench often waits for
    your module to do something, with a line like \code{wait until tx = '0';}. If your module never
    does, that wait never finishes, and without a stop time the simulation hangs: no output, no
    error, just a terminal sitting there waiting. 10\,ms of \emph{simulated} time is far more than
    any testbench here needs, and takes well under a second in real time.
  \item A passing run prints its ``all checks passed'' note and exits with 0. Every testbench in
    this book ends with that note, so if you do not see it the run did not finish: it either failed
    a check or hit the stop time.
\end{itemize}

\subsection{Checking an exercise}
\label{c2:sub:exercise}
Every VHDL exercise has its testbench in a directory of its own. This chapter's
Exercise~\ref{c2:ex:xyz}, for example, has
\file{lectures/L02/exercises/xyz_logic/xyz_logic_tb.vhd}. Write your module in the \textbf{same
directory}, with exactly the name the exercise asks for, and then run the three commands from
there:

\begin{shell}
cd lectures/L02/exercises/xyz_logic
ghdl -a --std=93 xyz_logic.vhd xyz_logic_tb.vhd
ghdl -e --std=93 xyz_logic_tb
ghdl -r --std=93 xyz_logic_tb --assert-level=error --stop-time=10ms
\end{shell}

Every testbench's file header lists exactly these commands.

The testbench connects to your module \textbf{positionally}, not by name, so your entity must
declare its ports in the \textbf{same order} as the exercise lists them. The names are yours to
choose, but use the given ones anyway, so that your module and the testbench's error messages speak
about the same signals.

A port declared in the wrong order fails in one of two ways, neither of which says ``wrong order'':
\begin{itemize}
  \item if the misplaced ports have \textbf{different} types or widths, analysis stops with
    \code{can't associate "sel" with port "sel"}, pointing at the testbench's \code{port map} line.
  \item if they have the \textbf{same} type, say several \code{std_logic} ports, it analyses and
    elaborates perfectly well, and then fails at run time as an ordinary assertion failure, which
    reads exactly like a logic bug in your design.
\end{itemize}

So when a testbench fails on the very first case it checks: reread your port order before you go
through your architecture. The same goes for \textbf{generics}: a module that carries them must
declare exactly the ones the exercise states, in order.

\subsection{When the design needs more than one file}
\label{c2:sub:multifile}
Some exercises build one module out of another, so the \code{ghdl -a} line names several files
instead of two. The elaboration and run steps are unchanged: they still take only the testbench's
\textbf{entity name}.

\begin{shell}
ghdl -a --std=93 <subcomponent>.vhd <module>.vhd <module>_tb.vhd
ghdl -e --std=93 <module>_tb
ghdl -r --std=93 <module>_tb --assert-level=error --stop-time=10ms
\end{shell}

\begin{itemize}
  \item \textbf{Order matters on the analysis line}: name a module before anything that
    instantiates it, so that GHDL has already seen the entity when it reaches the
    \code{entity work.<name>} referring to it. Subcomponents first, your top level after them, the
    testbench last.
  \item No exercise directory has a subcomponent handed out with it. Every module a design
    instantiates is one you wrote in an earlier exercise and copy in yourself; the exercise says
    which, and from where.
  \item The exercises that need this are this chapter's \code{hex_display} (which instantiates your
    \code{display}), \chapref{c4:ch}'s \code{led_toggle_sync2}, \chapref{c7:ch}'s \code{blinker}
    and \code{walking_led}, and all of \chapref{c8:ch}'s, where the receiver names five files and
    the transmitter four.
  \item Every testbench's file header lists the exact command for its own exercise.
\end{itemize}

\subsection{Reading the result}
\label{c2:sub:result}
A failed check looks like this:

\begin{console}
xyz_logic_tb.vhd:38:13:@10ns:(assertion failure): xyz_logic: wrong X with a='0' b='0' c='0' d='0', expected '0' but got '1'!
ghdl:error: assertion failed
\end{console}

Read it in order: \code{file:line} is where the assertion is, \code{@time} is when it fired, and
after that comes the message. It is the same shape as a compilation error from GHDL, so it maps
onto something you already know. No error means every case matched.

\subsection{A complete example}
\label{c2:sub:example}
A combinational testbench in the form every testbench in this book uses, here for
\chapref{c1:ch}'s \code{or_gate}:

\begin{vhdlcode}
library ieee;
use ieee.std_logic_1164.all;
use ieee.numeric_std.all;

entity or_gate_tb is
end entity;

architecture behaviour of or_gate_tb is
signal a, b, x: std_logic;
begin
    dut: entity work.or_gate
        port map(a, b, x);   -- positional: same order as or_gate's port clause.

    SIM_PROCESS: process is
    begin
        for i in 0 to 3 loop
            (a, b) <= std_logic_vector(to_unsigned(i, 2));   -- apply inputs
            wait for 10 ns;                                  -- let them settle
            assert x = (a or b)                              -- check AFTER the wait
                report "or_gate: wrong output with a=" & std_logic'image(a)
                     & " b=" & std_logic'image(b)
                     & ", expected " & std_logic'image(a or b)
                     & " but got " & std_logic'image(x) & "!"
                severity failure;
        end loop;
        report "or_gate: all checks passed!" severity note;
        wait;                                                -- end the simulation
    end process;
end architecture;
\end{vhdlcode}

Two rules keep a combinational testbench correct:
\begin{itemize}
  \item \textbf{Apply, wait, check, in that order.} A signal assignment only \emph{schedules} a
    value, so the inputs and the output are not updated until the \code{wait}. Checking before it
    tests the previous round's values.
  \item \textbf{End with \code{wait;}} so that the stimulus process suspends and the simulation
    stops by itself.
\end{itemize}

Clocked designs need a little more, a clock-generating process and a way to stop it, which is
exactly the machinery the testbenches handed out already contain. For the exercises you write only
the module.

\paragraph{Run one now}
Against this chapter's first worked example:

\begin{shell}
cd lectures/L02/gate_network
ghdl -a --std=93 gate_network.vhd gate_network_tb.vhd
ghdl -e --std=93 gate_network_tb
ghdl -r --std=93 gate_network_tb --assert-level=error --stop-time=10ms
\end{shell}
